\scanpage{140}

% L1–4, bracketed lemma
[ \uline{Lemme} 5 ($S$ affine.) Soit $F\in\Delta_{\{n\}}(E,q)$, où $N=\mathrm{rg}\,E$ \deleted{\unsure{$N$}} $=2n$. Alors $\exists$ supp. $G$ de $F$ tel que $q\mid G=0$. (La forme $q$ sur $E\simeq F\times G$ s'exprime donc par

% L5–6, display
\[
\begin{aligned}
q(x\oplus y)&=\varphi\deleted{(}\langle x,y\rangle &&(=\langle x,y\rangle\ \text{si}\\
&&&\text{on identifie }G\unsure{:|\text{ à }}F^\vee)
\end{aligned}
\]

% L7–8, proof, bracket continues
\uline{Dém} On choisit un suppl$^{\mathrm t}$ quelc, puis on le corrige, en utilisant que

% L9, display and end of bracket
\[
\mathrm{Hom}\bigl(\underbrace{E/F}_{\substack{\shortparallel\\ \deleted{F}G}},\underset{\substack{\shortparallel\\G^\vee}}{F}\bigr)
\longrightarrow\mathrm{Quad}(G)\quad\text{est surjectif}\ldots\ \bigr]
\]

% L10–12, middle text
Donc dans le raisonnement précédent, \unsure{ops|ppg|OPS} $q\mid G=0$, de sorte que le calcul du discriminant donne

% L13, display
\[
\delta(q)=(-1)^n(\omega_F)^2
\]

% L14, text
ou encore

% L15, boxed display
\[
\boxed{(\omega_F)^2=(-1)^n\delta(q)^{-1}}
\]

% L16, cancelled line
\deleted{Considérons d'autre part}

% L17, lower text
On définit par $F\longmapsto\omega_F$ un morphisme

% L18–21, display and annotations
\[
\begin{aligned}
\underline{\Delta}_{\{n\}}(E,q)&\longrightarrow
M((-1)^n\delta(q))\overset{\substack{\text{\unsure{déf|dit} des variétés}\\\text{d'Arf}\ldots}}{=}\underline{\mathrm{Arf}}(E,q)\\[-2pt]
&\hspace{3em}\shortparallel\\[-2pt]
&\quad\text{sous-schéma de }V(\underline{\det}\,E)\\
&\qquad\text{des racines carrées de}\\
&\qquad (-1)^n\delta(q)^{-1}
\end{aligned}
\]

% bottom left, pencil
% archive sheet number: 139
